\documentclass[11pt,a4paper]{article}
\usepackage{fontspec,amsmath,amssymb,amsthm,microtype}
\setmainfont{lmroman10-regular.otf}[BoldFont=lmroman10-bold.otf,ItalicFont=lmroman10-italic.otf,BoldItalicFont=lmroman10-bolditalic.otf]
\usepackage{xcolor}
\usepackage[most]{tcolorbox}
\usepackage{mathrsfs}
\providecommand{\lt}{<}
\providecommand{\gt}{>}
\setmonofont{DejaVuSansMono.ttf}[BoldFont=DejaVuSansMono-Bold.ttf,ItalicFont=DejaVuSansMono-Oblique.ttf,BoldItalicFont=DejaVuSansMono-BoldOblique.ttf,Scale=0.86]
\newfontfamily\leanfallback{STIXTwoMath-Regular.otf}[Scale=0.86]
\definecolor{leanrule}{gray}{0.62}
\definecolor{leanback}{gray}{0.962}
\definecolor{leantitle}{gray}{0.22}
\definecolor{leanmuted}{gray}{0.38}
\newtcolorbox{leanbox}[1]{enhanced,breakable,colback=leanback,colframe=leanrule,
  boxrule=0.4pt,arc=0.6mm,left=2.4mm,right=2.4mm,top=1.8mm,bottom=1.8mm,
  before skip=0.8em,after skip=1em,fonttitle=\small\bfseries,coltitle=leantitle,
  colbacktitle=leanback,titlerule=0.3pt,title={#1}}
\newlength{\leancharwidth}
\newenvironment{leancode}{\par\smallskip\begingroup\ttfamily\footnotesize\raggedright
  \settowidth{\leancharwidth}{M}\setlength{\parskip}{0pt}\setlength{\parindent}{0pt}
  \frenchspacing}{\par\endgroup\smallskip}
\newcommand{\leanline}[2]{\hangindent=\dimexpr#1\leancharwidth+4\leancharwidth\relax
  \hangafter=1\noindent\hspace*{#1\leancharwidth}#2\par}
\newcommand{\leansource}[1]{\par{\footnotesize\color{leanmuted}#1}\par}
\newcommand{\leanlabel}[1]{\par\smallskip\noindent{\small\bfseries #1}\par}
\usepackage[margin=25mm]{geometry}
\usepackage[unicode,hidelinks]{hyperref}
\hypersetup{pdfauthor={LangevinLean Project}}
\setlength{\emergencystretch}{3em}
\setlength{\parskip}{0.4em}
\setlength{\parindent}{0pt}
\allowdisplaybreaks[2]
\raggedbottom
\title{\texorpdfstring{The KLMC2 Gaussian Markov law and its coupling representation}{The KLMC2 Gaussian Markov law and its coupling representation}}
\author{LangevinLean Project}
\date{}
\begin{document}\maketitle
Draft mathematical exposition; these notes may contain mistakes.


The second-order sampler is now a single Gaussian Markov law, not a law depending on an analysis coupling. Its four-noise stream is built from the specified Brownian kernels. The module proves the exact complete initial/\allowbreak{}noise product law, identifies the original pointwise output with the correlated-Gaussian Markov iterates, and identifies every random-Hessian \(L^{2}\) coupling construction with that same output.



The specified Brownian source and original corrected update define one sampler law. Whole-stream independence, coupling representatives, and finite transport connect that law to the second-order error analysis.


\section{\texorpdfstring{The independent-source second-order law}{The independent-source second-order law}}


The source contains independent standard Gaussian velocity and a vector Brownian coordinate. The defined four-kernel noise reads that Brownian coordinate. Its finite second moments follow from the convolution construction and measure-preserving source projection. The pointwise original KLMC2 chain on initial-position times source defines the specified output law, which is a probability measure by the proved original-step measurability.



The initial position is independent of the standard Gaussian velocity and the Brownian source. The four-response endpoint map is measurable and has the constructed stream as its input. Its pushforward defines the constructed KLMC2 position law without reference to a coupling.

\begin{leanbox}{Lean: \texttt{SourceSample}, \texttt{sourceInnovation}, \texttt{sourceInnovation\_\allowbreak{}memLp} and 2 more}
\leansource{\texttt{SourceSample} \quad KineticLangevin/\allowbreak{}CanonicalKLMC2Law.lean, lines 11–11}
\begin{leancode}
\leanline{0}{abbrev\ SourceSample\ :=\ H\ ι\ ×\ ULA.CanonicalBrownian.VectorSample\ ι}
\end{leancode}
\leansource{\texttt{sourceInnovation} \quad KineticLangevin/\allowbreak{}CanonicalKLMC2Law.lean, lines 14–15}
\begin{leancode}
\leanline{0}{def\ sourceInnovation\ (γ\ :\ ℝ)\ (h\ :\ ℝ≥0)\ (n\ :\ ℕ)\ (p\ :\ SourceSample\ ι)\ :\ Fin\ 4\ →\ H\ ι\ :=}
\leanline{2}{innovation\ ι\ γ\ h\ n\ p.2}
\end{leancode}
\leansource{\texttt{sourceInnovation\_\allowbreak{}memLp} \quad KineticLangevin/\allowbreak{}CanonicalKLMC2Law.lean, lines 17–19}
\begin{leancode}
\leanline{0}{theorem\ sourceInnovation\_memLp\ (γ\ :\ ℝ)\ (h\ :\ ℝ≥0)\ (n\ :\ ℕ)\ (j\ :\ Fin\ 4)\ :}
\leanline{4}{MemLp\ (fun\ p\ =>\ sourceInnovation\ ι\ γ\ h\ n\ p\ j)\ 2\ (GenericCanonicalSource.canonicalSourceProbability\ ι)}
\end{leancode}
\leansource{\texttt{law} \quad KineticLangevin/\allowbreak{}CanonicalKLMC2Law.lean, lines 22–25}
\begin{leancode}
\leanline{0}{def\ law\ (g\ :\ H\ ι\ →\ H\ ι)\ (Hess\ :\ H\ ι\ →\ H\ ι\ →L[ℝ]\ H\ ι)\ (γ\ :\ ℝ)\ (h\ :\ ℝ≥0)}
\leanline{4}{(μ\ :\ Measure\ (H\ ι))\ (n\ :\ ℕ)\ :\ Measure\ (H\ ι)\ :=}
\leanline{2}{productKLMC2Law\ (GenericCanonicalSource.canonicalSourceProbability\ ι)\ g\ Hess\ γ\ h\ μ\ Prod.fst}
\leanline{4}{(fun\ k\ =>\ sourceInnovation\ ι\ γ\ h\ k)\ n}
\end{leancode}
\leansource{\texttt{law\_\allowbreak{}probability} \quad KineticLangevin/\allowbreak{}CanonicalKLMC2Law.lean, lines 27–32}
\begin{leancode}
\leanline{0}{theorem\ law\_probability\ (g\ :\ H\ ι\ →\ H\ ι)\ (Hess\ :\ H\ ι\ →\ H\ ι\ →L[ℝ]\ H\ ι)}
\leanline{4}{(hG\ :\ Continuous\ g)\ (hH\ :\ Continuous\ Hess)\ (γ\ :\ ℝ)\ (h\ :\ ℝ≥0)}
\leanline{4}{(μ\ :\ Measure\ (H\ ι))\ [IsProbabilityMeasure\ μ]\ (n\ :\ ℕ)\ :\ IsProbabilityMeasure\ (law\ ι\ g\ Hess\ γ\ h\ μ\ n)}
\end{leancode}
\leanlabel{How the Lean proof follows the argument}
\texttt{Source\allowbreak{}Sample} names the mathematical object defined above. Its defining expression appears in the adjacent statement; subsequent identities may unfold this definition directly.




\texttt{source\allowbreak{}Innovation} names the mathematical object defined above. Its defining expression appears in the adjacent statement; subsequent identities may unfold this definition directly.




\texttt{source\allowbreak{}Innovation\_\allowbreak{}mem\allowbreak{}Lp} gives the formal counterpart of this step.




\texttt{law} names the mathematical object defined above. Its defining expression appears in the adjacent statement; subsequent identities may unfold this definition directly.




\texttt{law\_\allowbreak{}probability} gives the formal counterpart of this step.
\end{leanbox}


Earlier results used here: \cite[Section~3.8.16]{LangevinLean2026ULASC}.

\section{\texorpdfstring{The complete Gaussian Markov identification}{The complete Gaussian Markov identification}}


On phase-times-Brownian source, the initial phase is independent of the entire iid four-kernel stream. Reassociation of independent initial position and standard Gaussian velocity pulls this exact law back to the sampler’s source. The resulting initial-phase/\allowbreak{}whole-stream product law determines the Markov iterates through the general random-recursion theorem. Projecting to position proves the independently defined specified Brownian sampler law equals exactly the defined correlated-Gaussian Markov output.



The whole-stream product law implies that every next response block is independent of the history generated by the recursion. This is the precise hypothesis needed for iterated Markov transitions. Pushing the full input law through the recursion identifies the constructed and Markov descriptions.

\begin{leanbox}{Lean: \texttt{initial\_\allowbreak{}innovation\_\allowbreak{}jointLaw}, \texttt{law\_\allowbreak{}eq\_\allowbreak{}markovLaws}}
\leansource{\texttt{initial\_\allowbreak{}innovation\_\allowbreak{}jointLaw} \quad KineticLangevin/\allowbreak{}CanonicalKLMC2Law.lean, lines 35–43}
\begin{leancode}
\leanline{0}{theorem\ initial\_innovation\_jointLaw\ (μ\ :\ Measure\ (H\ ι))\ [IsProbabilityMeasure\ μ]\ (γ\ :\ ℝ)\ (h\ :\ ℝ≥0)\ :}
\leanline{4}{HasLaw\ (fun\ p\ :\ H\ ι\ ×\ SourceSample\ ι\ =>\ ((p.2.1,p.1),fun\ n\ =>\ innovation\ ι\ γ\ h\ n\ p.2.2))}
\leanline{6}{(((stdGaussian\ (H\ ι)).prod\ μ).prod\ (Measure.infinitePi\ (fun\ \_\ :\ ℕ\ =>\ gaussianInnovationLaw\ ι\ γ\ h)))}
\leanline{6}{(μ.prod\ (GenericCanonicalSource.canonicalSourceProbability\ ι))}
\end{leancode}
\leansource{\texttt{law\_\allowbreak{}eq\_\allowbreak{}markovLaws} \quad KineticLangevin/\allowbreak{}CanonicalKLMC2Law.lean, lines 46–69}
\begin{leancode}
\leanline{0}{theorem\ law\_eq\_markovLaws\ (g\ :\ H\ ι\ →\ H\ ι)\ (Hess\ :\ H\ ι\ →\ H\ ι\ →L[ℝ]\ H\ ι)}
\leanline{4}{(hG\ :\ Continuous\ g)\ (hH\ :\ Continuous\ Hess)\ (γ\ :\ ℝ)\ (h\ :\ ℝ≥0)}
\leanline{4}{(μ\ :\ Measure\ (H\ ι))\ [IsProbabilityMeasure\ μ]\ (n\ :\ ℕ)\ :}
\leanline{4}{law\ ι\ g\ Hess\ γ\ h\ μ\ n\ =}
\leanline{6}{(ULA.markovLaws\ (klmc2Step\ g\ Hess\ γ\ h)\ (gaussianInnovationLaw\ ι\ γ\ h)}
\leanline{8}{((stdGaussian\ (H\ ι)).prod\ μ)\ n).map\ Prod.snd}
\end{leancode}
\leanlabel{How the Lean proof follows the argument}
\texttt{initial\_\allowbreak{}innovation\_\allowbreak{}joint\allowbreak{}Law} gives the formal counterpart of this step. The main intermediate facts are \texttt{hj} in that order. Named proof applications include \texttt{hj.\allowbreak{}comp}; the adjacent proof shows their hypotheses and order of use.




\texttt{law\_\allowbreak{}eq\_\allowbreak{}markov\allowbreak{}Laws} gives the formal counterpart of this step. The main intermediate facts are \texttt{P}, \texttt{hs}, \texttt{hc}, \texttt{hl}, \texttt{hp} in that order. Named proof applications include \texttt{hp}; the adjacent proof shows their hypotheses and order of use.
\end{leanbox}


Earlier results used here: \cite[Section~3.17.1]{LangevinLean2026ULASC}; \cite[Section~3.20.1]{LangevinLean2026ULASC}; \cite[Section~3.20.4]{LangevinLean2026ULASC}; \cite[Section~3.33.6]{LangevinLean2026ULASC}; \cite[Section~3.8.17]{LangevinLean2026ULASC}.

\section{\texorpdfstring{Identify every \(L^{2}\) coupling construction}{Identify every L² coupling construction}}


The directly lifted Brownian convolutions and the generic independent-source noise extension have the same original source representatives, hence are equal in \(L^{2}\) coordinate by coordinate. The numerical chain uses the original gradient lift, the bounded random Hessian and those same four convolutions. The general coupling-independent output theorem therefore gives its exact specified law. The zero and successor identities give the initial state and recurrence needed by the stochastic convergence proof.



The direct Brownian response lift and the generic source lift agree almost everywhere in each coordinate. With the same original gradient and random Hessian, their numerical updates coincide. This connects the constructed coupling variables to the single already defined sampler law.

\begin{leanbox}{Lean: \texttt{extensionInnovation\_\allowbreak{}eq\_\allowbreak{}generic}, \texttt{numericalChain}, \texttt{numericalChain\_\allowbreak{}outputLaw} and 2 more}
\leansource{\texttt{extensionInnovation\_\allowbreak{}eq\_\allowbreak{}generic} \quad KineticLangevin/\allowbreak{}CanonicalKLMC2Law.lean, lines 74–83}
\begin{leancode}
\leanline{0}{theorem\ extensionInnovation\_eq\_generic\ (c\ :\ ULA.Coupling\ μ\ π)\ (γ\ :\ ℝ)\ (h\ :\ ℝ≥0)\ (n\ :\ ℕ)\ :}
\leanline{4}{extensionInnovation\ c\ γ\ h\ n\ =}
\leanline{6}{extensionKLMC2Noise\ (GenericCanonicalSource.canonicalSourceProbability\ ι)\ c}
\leanline{8}{(fun\ k\ =>\ sourceInnovation\ ι\ γ\ h\ k)\ (sourceInnovation\_memLp\ ι\ γ\ h)\ n}
\end{leancode}
\leansource{\texttt{numericalChain} \quad KineticLangevin/\allowbreak{}CanonicalKLMC2Law.lean, lines 86–94}
\begin{leancode}
\leanline{0}{def\ numericalChain\ (c\ :\ ULA.Coupling\ μ\ π)\ (hc\ :\ ULA.couplingCost\ c.joint\ ≠\ ⊤)}
\leanline{4}{(hπ\ :\ MemLp\ (id\ :\ H\ ι\ →\ H\ ι)\ 2\ π)}
\leanline{4}{(g\ :\ H\ ι\ →\ H\ ι)\ (Hess\ :\ H\ ι\ →\ H\ ι\ →L[ℝ]\ H\ ι)\ (M\ M2\ :\ ℝ≥0)}
\leanline{4}{(hG\ :\ LipschitzWith\ M\ g)\ (hH\ :\ LipschitzWith\ M2\ Hess)\ (hg\ :\ ∀\ x,HasFDerivAt\ g\ (Hess\ x)\ x)}
\leanline{4}{(γ\ :\ ℝ)\ (h\ :\ ℝ≥0)\ :\ ℕ\ →\ Lp\ (H\ ι)\ 2\ (GenericCanonicalCoupling.probability\ c)\ ×}
\leanline{6}{Lp\ (H\ ι)\ 2\ (GenericCanonicalCoupling.probability\ c)\ :=}
\leanline{2}{klmc2Chain\ (ULA.lpDrift\ (GenericCanonicalCoupling.probability\ c)\ g\ hG)}
\leanline{4}{(lpHessian\ (GenericCanonicalCoupling.probability\ c)\ g\ Hess\ M\ M2\ hG\ hH\ hg)\ γ\ h}
\leanline{4}{(GenericCanonicalCoupling.numericalInitial\ c\ hc\ hπ)\ (extensionInnovation\ c\ γ\ h)}
\end{leancode}
\leansource{\texttt{numericalChain\_\allowbreak{}outputLaw} \quad KineticLangevin/\allowbreak{}CanonicalKLMC2Law.lean, lines 97–109}
\begin{leancode}
\leanline{0}{theorem\ numericalChain\_outputLaw\ (c\ :\ ULA.Coupling\ μ\ π)\ (hc\ :\ ULA.couplingCost\ c.joint\ ≠\ ⊤)}
\leanline{4}{(hπ\ :\ MemLp\ (id\ :\ H\ ι\ →\ H\ ι)\ 2\ π)}
\leanline{4}{(g\ :\ H\ ι\ →\ H\ ι)\ (Hess\ :\ H\ ι\ →\ H\ ι\ →L[ℝ]\ H\ ι)\ (M\ M2\ :\ ℝ≥0)}
\leanline{4}{(hG\ :\ LipschitzWith\ M\ g)\ (hH\ :\ LipschitzWith\ M2\ Hess)\ (hg\ :\ ∀\ x,HasFDerivAt\ g\ (Hess\ x)\ x)}
\leanline{4}{(γ\ :\ ℝ)\ (h\ :\ ℝ≥0)\ (n\ :\ ℕ)\ :}
\leanline{4}{(GenericCanonicalCoupling.probability\ c).map\ (numericalChain\ c\ hc\ hπ\ g\ Hess\ M\ M2\ hG\ hH\ hg\ γ\ h\ n).2\ =}
\leanline{6}{law\ ι\ g\ Hess\ γ\ h\ μ\ n}
\end{leancode}
\leansource{\texttt{numericalChain\_\allowbreak{}zero} \quad KineticLangevin/\allowbreak{}CanonicalKLMC2Law.lean, lines 112–117}
\begin{leancode}
\leanline{0}{@[simp]\ theorem\ numericalChain\_zero\ (c\ :\ ULA.Coupling\ μ\ π)\ (hc\ :\ ULA.couplingCost\ c.joint\ ≠\ ⊤)}
\leanline{4}{(hπ\ :\ MemLp\ (id\ :\ H\ ι\ →\ H\ ι)\ 2\ π)}
\leanline{4}{(g\ :\ H\ ι\ →\ H\ ι)\ (Hess\ :\ H\ ι\ →\ H\ ι\ →L[ℝ]\ H\ ι)\ (M\ M2\ :\ ℝ≥0)}
\leanline{4}{(hG\ :\ LipschitzWith\ M\ g)\ (hH\ :\ LipschitzWith\ M2\ Hess)\ (hg\ :\ ∀\ x,HasFDerivAt\ g\ (Hess\ x)\ x)}
\leanline{4}{(γ\ :\ ℝ)\ (h\ :\ ℝ≥0)\ :}
\leanline{4}{numericalChain\ c\ hc\ hπ\ g\ Hess\ M\ M2\ hG\ hH\ hg\ γ\ h\ 0=GenericCanonicalCoupling.numericalInitial\ c\ hc\ hπ}
\end{leancode}
\leansource{\texttt{numericalChain\_\allowbreak{}succ} \quad KineticLangevin/\allowbreak{}CanonicalKLMC2Law.lean, lines 120–128}
\begin{leancode}
\leanline{0}{theorem\ numericalChain\_succ\ (c\ :\ ULA.Coupling\ μ\ π)\ (hc\ :\ ULA.couplingCost\ c.joint\ ≠\ ⊤)}
\leanline{4}{(hπ\ :\ MemLp\ (id\ :\ H\ ι\ →\ H\ ι)\ 2\ π)}
\leanline{4}{(g\ :\ H\ ι\ →\ H\ ι)\ (Hess\ :\ H\ ι\ →\ H\ ι\ →L[ℝ]\ H\ ι)\ (M\ M2\ :\ ℝ≥0)}
\leanline{4}{(hG\ :\ LipschitzWith\ M\ g)\ (hH\ :\ LipschitzWith\ M2\ Hess)\ (hg\ :\ ∀\ x,HasFDerivAt\ g\ (Hess\ x)\ x)}
\leanline{4}{(γ\ :\ ℝ)\ (h\ :\ ℝ≥0)\ (n\ :\ ℕ)\ :}
\leanline{4}{numericalChain\ c\ hc\ hπ\ g\ Hess\ M\ M2\ hG\ hH\ hg\ γ\ h\ (n+1)=}
\leanline{6}{klmc2Step\ (ULA.lpDrift\ (GenericCanonicalCoupling.probability\ c)\ g\ hG)}
\leanline{8}{(lpHessian\ (GenericCanonicalCoupling.probability\ c)\ g\ Hess\ M\ M2\ hG\ hH\ hg)\ γ\ h}
\leanline{8}{(numericalChain\ c\ hc\ hπ\ g\ Hess\ M\ M2\ hG\ hH\ hg\ γ\ h\ n)\ (extensionInnovation\ c\ γ\ h\ n)}
\end{leancode}
\leanlabel{How the Lean proof follows the argument}
\texttt{extension\allowbreak{}Innovation\_\allowbreak{}eq\_\allowbreak{}generic} gives the formal counterpart of this step. Named proof applications include \texttt{Lp.\allowbreak{}ext}; the adjacent proof shows their hypotheses and order of use.




\texttt{numerical\allowbreak{}Chain} names the mathematical object defined above. Its defining expression appears in the adjacent statement; subsequent identities may unfold this definition directly.




\texttt{numerical\allowbreak{}Chain\_\allowbreak{}output\allowbreak{}Law} gives the formal counterpart of this step. The main intermediate facts are \texttt{h\allowbreak{}N} in that order. Named proof applications include \texttt{product\_\allowbreak{}klmc2\allowbreak{}Chain\_\allowbreak{}output\allowbreak{}Law}; the adjacent proof shows their hypotheses and order of use.




\texttt{numerical\allowbreak{}Chain\_\allowbreak{}zero} is a definitional equality: \texttt{rfl} closes it because both sides reduce to the same expression.




\texttt{numerical\allowbreak{}Chain\_\allowbreak{}succ} is a definitional equality: \texttt{rfl} closes it because both sides reduce to the same expression.
\end{leanbox}


Earlier results used here: \cite[Section~3.13.2]{LangevinLean2026ULASC}; \cite[Section~3.17.1]{LangevinLean2026ULASC}; \cite[Section~3.17.2]{LangevinLean2026ULASC}.

\section{\texorpdfstring{Establish finite transport and the exact initial cost}{Establish finite transport and the exact initial cost}}


A finite-cost initial coupling and the finite target second moment construct an \(L^{2}\) numerical chain. Its output-law equality transfers this finite second moment to the single constructed sampler law. Existence of a finite-cost initial plan therefore proves the output Wasserstein distance is finite. Finally, the initial numerical/\allowbreak{}reference position discrepancy is exactly the root-mean-square cost of the supplied initial coupling. These facts permit the final convergence estimate to take the infimum over initial couplings without assuming an optimizer or hiding infinite distances through real conversion.



The finite Gibbs moment and one finite-cost initial plan give a square-integrable output coupling. This proves finite output Wasserstein distance before taking real values. The initial discrepancy norm equals the chosen plan cost, enabling the final infimum over all such plans.

\begin{leanbox}{Lean: \texttt{law\_\allowbreak{}memLp\_\allowbreak{}of\_\allowbreak{}finiteCost}, \texttt{law\_\allowbreak{}wasserstein2\_\allowbreak{}ne\_\allowbreak{}top}, \texttt{initial\_\allowbreak{}position\_\allowbreak{}cost}}
\leansource{\texttt{law\_\allowbreak{}memLp\_\allowbreak{}of\_\allowbreak{}finiteCost} \quad KineticLangevin/\allowbreak{}CanonicalKLMC2Law.lean, lines 131–138}
\begin{leancode}
\leanline{0}{theorem\ law\_memLp\_of\_finiteCost\ (c\ :\ ULA.Coupling\ μ\ π)\ (hc\ :\ ULA.couplingCost\ c.joint\ ≠\ ⊤)}
\leanline{4}{(hπ\ :\ MemLp\ (id\ :\ H\ ι\ →\ H\ ι)\ 2\ π)}
\leanline{4}{(g\ :\ H\ ι\ →\ H\ ι)\ (Hess\ :\ H\ ι\ →\ H\ ι\ →L[ℝ]\ H\ ι)\ (M\ M2\ :\ ℝ≥0)}
\leanline{4}{(hG\ :\ LipschitzWith\ M\ g)\ (hH\ :\ LipschitzWith\ M2\ Hess)\ (hg\ :\ ∀\ x,HasFDerivAt\ g\ (Hess\ x)\ x)}
\leanline{4}{(γ\ :\ ℝ)\ (h\ :\ ℝ≥0)\ (n\ :\ ℕ)\ :}
\leanline{4}{MemLp\ (id\ :\ H\ ι\ →\ H\ ι)\ 2\ (law\ ι\ g\ Hess\ γ\ h\ μ\ n)}
\end{leancode}
\leansource{\texttt{law\_\allowbreak{}wasserstein2\_\allowbreak{}ne\_\allowbreak{}top} \quad KineticLangevin/\allowbreak{}CanonicalKLMC2Law.lean, lines 141–150}
\begin{leancode}
\leanline{0}{theorem\ law\_wasserstein2\_ne\_top\ [IsProbabilityMeasure\ π]}
\leanline{4}{(hπ\ :\ MemLp\ (id\ :\ H\ ι\ →\ H\ ι)\ 2\ π)}
\leanline{4}{(g\ :\ H\ ι\ →\ H\ ι)\ (Hess\ :\ H\ ι\ →\ H\ ι\ →L[ℝ]\ H\ ι)\ (M\ M2\ :\ ℝ≥0)}
\leanline{4}{(hG\ :\ LipschitzWith\ M\ g)\ (hH\ :\ LipschitzWith\ M2\ Hess)\ (hg\ :\ ∀\ x,HasFDerivAt\ g\ (Hess\ x)\ x)}
\leanline{4}{(γ\ :\ ℝ)\ (h\ :\ ℝ≥0)\ (hfinite\ :\ ULA.wasserstein2\ μ\ π\ ≠\ ⊤)\ (n\ :\ ℕ)\ :}
\leanline{4}{ULA.wasserstein2\ (law\ ι\ g\ Hess\ γ\ h\ μ\ n)\ π\ ≠\ ⊤}
\end{leancode}
\leansource{\texttt{initial\_\allowbreak{}position\_\allowbreak{}cost} \quad KineticLangevin/\allowbreak{}CanonicalKLMC2Law.lean, lines 153–160}
\begin{leancode}
\leanline{0}{theorem\ initial\_position\_cost\ (c\ :\ ULA.Coupling\ μ\ π)\ (hc\ :\ ULA.couplingCost\ c.joint\ ≠\ ⊤)}
\leanline{4}{(hπ\ :\ MemLp\ (id\ :\ H\ ι\ →\ H\ ι)\ 2\ π)}
\leanline{4}{(g\ :\ H\ ι\ →\ H\ ι)\ (Hess\ :\ H\ ι\ →\ H\ ι\ →L[ℝ]\ H\ ι)\ (M\ M2\ :\ ℝ≥0)}
\leanline{4}{(hG\ :\ LipschitzWith\ M\ g)\ (hH\ :\ LipschitzWith\ M2\ Hess)\ (hg\ :\ ∀\ x,HasFDerivAt\ g\ (Hess\ x)\ x)}
\leanline{4}{(γ\ :\ ℝ)\ (h\ :\ ℝ≥0)\ :}
\leanline{4}{‖(numericalChain\ c\ hc\ hπ\ g\ Hess\ M\ M2\ hG\ hH\ hg\ γ\ h\ 0).2-(GenericCanonicalCoupling.referenceInitial\ c\ hπ).2‖=}
\leanline{6}{(ULA.couplingCost\ c.joint).toReal}
\end{leancode}
\leanlabel{How the Lean proof follows the argument}
\texttt{law\_\allowbreak{}mem\allowbreak{}Lp\_\allowbreak{}of\_\allowbreak{}finite\allowbreak{}Cost} gives the formal counterpart of this step.




\texttt{law\_\allowbreak{}wasserstein2\_\allowbreak{}ne\_\allowbreak{}top} gives the formal counterpart of this step. Named proof applications include \texttt{ULA.\allowbreak{}wasserstein2\_\allowbreak{}ne\_\allowbreak{}top\_\allowbreak{}of\_\allowbreak{}mem\allowbreak{}Lp}; the adjacent proof shows their hypotheses and order of use.




\texttt{initial\_\allowbreak{}position\_\allowbreak{}cost} gives the formal counterpart of this step.
\end{leanbox}


Earlier results used here: \cite[Section~3.17.1]{LangevinLean2026ULASC}; \cite[Section~3.17.2]{LangevinLean2026ULASC}; \cite[Section~3.17.3]{LangevinLean2026ULASC}; \cite[Section~3.18.2]{LangevinLean2026ULASC}; \cite[Section~3.21.2]{LangevinLean2026ULASC}; \cite[Section~3.22.23]{LangevinLean2026ULASC}.

\begin{thebibliography}{ULA-LSI}
\bibitem[ULA-SC]{LangevinLean2026ULASC}
LangevinLean Project. \emph{Unadjusted Langevin under Strong Convexity}. LL-01, 2026. Draft report, content revision d34ead06987c. \url{https://odakyildiz.github.io/LangevinLean/docs/pdf/UnadjustedLangevin/FullProof.pdf}.
\end{thebibliography}

\end{document}
